% TOP_PROBLEM: {"id": 10000073, "title": "Absence of bigeodesics in first-passage percolation", "category_id": 19, "category": {"id": 19, "name": "probability", "display_name": "Probability", "description": "Problems involving probability theory, stochastic processes, and random structures.", "slug": "probability", "order_index": 19, "created_at": "2026-07-31T15:26:25.671Z"}, "status": "partially_solved"}
% FIRST_POSTED: 2026-10-06
% ATTEMPT_STATUS: unresolved
% OPENAI_CLAIM: {"schema_version":1,"source_repo":"https://github.com/openai/math","source_commit":"adc7f1241b42e322a6451854ab7e4b4c146bf78a","release_date":"2026-10-06","match":"partial","resolution":"partial","independently_reviewed":false,"second_pass":"n/a","families":[{"family":"212","title":"Planar first-passage geometry and the absence of bigeodesics","subject":"Probability and statistical mechanics","manuscripts":[{"title":"No bigeodesics in planar first-passage percolation","pdf_path":"preprints/No-bigeodesics-in-planar-first-passage-percolation-September-24-2026/main.pdf","readme_path":"preprints/No-bigeodesics-in-planar-first-passage-percolation-September-24-2026/README.md","date":"2026-09-24","theorem_ref":"Theorem (label thm:main), preprints/No-bigeodesics-in-planar-first-passage-percolation-September-24-2026/build/sections/01-introduction.tex:24"}],"lean":{"doc_path":"lean/docs/212.md","comparators":[],"declarations":[],"covers_main_theorem":false},"reasoning_trace":null}]}
\documentclass[11pt]{article}
\usepackage[margin=25mm]{geometry}
\usepackage{fontspec}
\setmainfont{Latin Modern Roman}
\usepackage{amsmath,amssymb,mathtools}
\usepackage{unicode-math}
\setmathfont{Latin Modern Math}
\usepackage{enumitem}
\usepackage{xurl,hyperref,bookmark}
\hypersetup{colorlinks=true,urlcolor=blue}
\setcounter{secnumdepth}{0}
\setlength{\parindent}{0pt}
\setlength{\parskip}{5pt}
\setlength{\emergencystretch}{3em}
\begin{document}
\section*{Absence of bigeodesics in first-passage percolation}

\subsection{OpenAI's claim}
OpenAI presents a mathematical result relevant to this question in the selected manuscript. Planar iid nonatomic first-passage weights are claimed to admit no doubly infinite minimizing path when the minimum of four weights has a finite second moment.

\subsection{Outline of the argument}
Busemann labels reduce potential bigeodesics to compact sectors. Finite increases of passage times create marked paths with a definite extra cost; deterministic label bins and telescoping Busemann comparisons impose a finite budget that cannot support all the marks, including environment-selected exceptional directions.

\subsection{Scope relative to this problem}
The theorem excludes every two-sided geodesic in the planar nearest-neighbor model under a continuous edge law and the minimum-of-four second-moment assumption. The problem statement asks natural iid models on Z\textasciicircum{}d without limiting d to two or imposing those law hypotheses. This entry records partial support; complete coverage of the question is not established by this catalogue.

\subsection{Formal verification}
For family 212, the documented entries give differentiability and continuously differentiable boundaries for planar Gamma shapes, including the exponential law. They do not supply strict convexity or the separate no-bigeodesic theorem. No Comparator is assigned to the selected manuscript here; the family page selects other papers or supporting statements. This site has not executed the Lean code. The scope descriptions do not certify the remaining manuscript arguments.

\subsection{Status caveat}
OpenAI attributes these manuscripts to work by an internal model. Its release cautions that arguments without formalization can contain mistakes. This catalogue has not independently checked the proofs; the labels report the mathematical claims attributed to OpenAI.

\subsection{Sources}
\begin{itemize}
\item OpenAI, \emph{No bigeodesics in planar first-passage percolation}, 2026-09-24; citation key \texttt{OAI:No-bigeodesics-in-planar-first-passage-percolation-September-24-2026}. \href{https://github.com/openai/math/blob/adc7f1241b42e322a6451854ab7e4b4c146bf78a/preprints/No-bigeodesics-in-planar-first-passage-percolation-September-24-2026/main.pdf}{PDF (GitHub)}; \href{https://github.com/openai/math/raw/adc7f1241b42e322a6451854ab7e4b4c146bf78a/preprints/No-bigeodesics-in-planar-first-passage-percolation-September-24-2026/main.pdf}{PDF (direct)}; \href{https://github.com/openai/math/blob/adc7f1241b42e322a6451854ab7e4b4c146bf78a/preprints/No-bigeodesics-in-planar-first-passage-percolation-September-24-2026/README.md}{Manuscript page and citation}; \href{https://github.com/openai/math/blob/main/preprints/No-bigeodesics-in-planar-first-passage-percolation-September-24-2026/main.pdf}{Latest version}.
\item Family 212: \href{https://github.com/openai/math/blob/adc7f1241b42e322a6451854ab7e4b4c146bf78a/lean/docs/212.md}{Lean scope}.
\item \href{https://github.com/openai/math/blob/adc7f1241b42e322a6451854ab7e4b4c146bf78a/CONTENTS.md}{OpenAI catalogue entry}.
\end{itemize}
\end{document}
